\documentclass{report}
\usepackage[utf8]{inputenc}
\usepackage[T1]{fontenc}
\usepackage{geometry}
\usepackage{graphicx}
\usepackage{float}
\usepackage{hyperref}
\usepackage[english,indonesian]{babel}

\title{Optimasi Pemrosesan Template LaTeX Menggunakan Pendekatan Berbasis AI}
\author{Academic Support Agent}
\date{\today}

\begin{document}
\maketitle

\begin{abstract}

\end{abstract}

\chapter{PENDAHULUAN}


\chapter{TINJAUAN PUSTAKA}


% Original content (commented out):
% %==================================================================
% Ini adalah file utama dari dokumen LaTeX ini
%==================================================================

%% DILARANG EDIT BAGIAN INI
\documentclass[12pt, a4paper, onecolumn, oneside, final]{report}

\input{identitas}
\input{preambles}
\input{database.hyphenate}

\begin{document}

\include{sampul}
\pagenumbering{roman}
\include{sampuldalam}
\include{abstrak}
\include{abstract}
\include{pernyataan}
\include{persetujuan}
\include{pengesahan}
\include{persembahan}
\include{katapengantar}
\include{daftar}
\pagenumbering{arabic}
\include{bab1}
\include{bab2}
\include{bab3}
\include{bab4}
\include{bab5}
\input{daftarpustaka}
\input{lampiran}

\end{document}
%% DILARANG EDIT BAGIAN INI

\end{document}